\chapter{はじめに}
\label{c-intro}

最短経路の原理は、人生と科学におけるほとんどの意思決定を導く：ある商品、人、あるいは1ビットの情報がある地点に存在し、目標地点に送る必要がある場合、可能な限り少ない労力で行うことが好まれる。これは通常、平面上では直線に沿って、より複雑な距離空間では測地線に沿ってアイテムを移動させることで達成される。Optimal transport の理論は、一度に1つのアイテムだけを移動させるのではなく、複数のアイテム（またはその連続分布）を同時にある配置から別の配置へ移動させる問題に関わる場合に、この直観を一般化するものである。学校教師が証言するかもしれないように、個人の集団の輸送を計画し、到着時に所定の目標配置に到達するという制約を課すことは、単一の個人に対して行うよりもはるかに複雑である。実際、集団や分布の観点で考えるには、\citet{Monge1781}の先駆的な研究で初めて示唆された、より高度な数学的形式化が必要である。しかし、その形式化が一見どれほど複雑に見えようとも、この問題は我々の日常生活と深く具体的な関係を持っている。人、商品、情報のいずれであれ、輸送が単一のアイテムの移動だけを含むことはきわめて稀である。ロジスティクス、生産計画、ネットワークルーティングにおけるすべての主要な経済問題は分布の移動を含み、そのテーマは optimal transport に関するすべての先駆的文献に現れる。実際、\citet{tolstoi1930methods}、\citet{Hitchcock41}、\citet{Kantorovich42}はすべて実用的な関心に導かれていた。数学者たちが、\citet{Brenier91}らの研究のおかげで、この理論が凸性、偏微分方程式、統計学との深い関係を持つ豊かな研究の場を提供することを発見したのは、それからわずか数年後、主に1980年代以降のことであった。2000年前後には、コンピュータ、イメージング、そしてより一般的にデータサイエンスの研究者たちが、optimal transport 理論が、bag-of-features や特徴記述子の形で得られる分布を比較するという、異なるより抽象的な文脈において分布を研究するための非常に強力なツールを提供することを理解した。

Optimal transport に関するいくつかの参考書が執筆されており、\citeauthor{Villani03}（\citeyear{Villani03,Villani09}）による最近の2つのモノグラフ、\citeauthor{rachev1998mass}（\citeyear{rachev1998mass,rachev1998mass2}）によるもの、そしてより最近では\citet{SantambrogioBook}によるものがある。これらの書籍に例示されるように、この理論におけるより形式的で抽象的な概念は、それ自体で数百ページに値する。Optimal transport が応用ツールとして徐々に確立されてきた今（例えば、\citet{galichon2016optimal}が最近提唱した経済学への応用）、我々はその豊かな文献と、データサイエンス、特にイメージングサイエンスと機械学習への応用を中心とした計算的観点とのバランスを取ることを試みた。この意味で、\citet{kolouri2017optimal}による最近のレビューの動機に従いつつ、より広い範囲をカバーすることを目指している。
%
究極的には、我々の目標は、OT の実用的有効性を支える主要な理論的知見の概観を提示し、これらの知見を高速な計算スキームに転換する方法の説明により多くの時間を費やすことである。
%
第\ref{c-continuous}章、第\ref{c-algo-basics}章、第\ref{c-entropic}章、第\ref{c-variational}章、第\ref{c-extensions}章の本文は、確率ベクトルまたは離散ヒストグラムの空間において optimal transport が誘導する幾何学の研究にのみ充てられている。
%
より上級の読者を対象として、同じ章の中で薄い灰色のボックスにおいて、離散測度向けに調整された optimal transport のより一般的な数学的説明も与える。離散測度は確率重みだけでなく、それらの重みが定義される位置によっても定義される。これらの位置は通常、連続距離空間に取られ、ランダム現象をモデル化するための第2の重要な自由度を与える。
%
最後に、第3の最も技術的な層の説明は濃い灰色のボックスで示され、離散である必要のない任意の測度を扱い、特に基底測度に関する密度を持ちうる。これは伝統的に OT 理論の古典的な教科書のほとんどにおけるデフォルトの設定であるが、一般的に実用的応用においてはあまり重要でない役割を果たす。
%
第\ref{c-algo-semidiscr}章から第\ref{c-statistical}章は連続測度と離散測度の相互作用を扱うため、より数学的な素養を持つ読者を対象としている。

Computational optimal transport の分野は、本稿執筆時点においてなお極めて活発な分野である。そのため、本サーベイでは触れていない多種多様なトピックが存在する。特に順序なく挙げると、distributionally robust optimization \citep{NIPS2015_5745,esfahani2018data,NIPS2018_7534,NIPS2018_8015}（パラメータ推定を、入力データの一定の Wasserstein 距離内の任意のデータ測度の最悪経験的リスクを最小化することによって行う手法）、Wasserstein 幾何学における Langevin Monte Carlo サンプリングアルゴリズムの収束 \citep{dalalyan2017user,pmlr-v65-dalalyan17a,pmlr-v75-bernton18a}、低次元設定における二乗ユークリッドコストを持つ OT を Monge-Amp\`ere 方程式を用いて解く他の数値手法~\citep{froese2011convergent,benamou2014numerical,sulman2011efficient}（注意~\ref{rem:MA}で簡単にのみ言及）がある。


%%%%%%%%%%%%%%%%%%%%%%%%%%%%%%%%%%
\section*{記法}

\begin{itemize}
	\item $\range{n}$：整数の集合 $\{1,\dots,n\}$。
	\item $\ones_{n,m}$：すべての成分が $1$ である $\RR^{n\times m}$ の行列。$\ones_{n}$：全成分が $1$ のベクトル。
	\item $\Identity_n$：サイズ $n\times n$ の単位行列。
	\item $u\in\RR^n$ に対し、$\diag(u)$ は対角成分が $u$ でそれ以外が零の $n\times n$ 行列。
	\item $\simplex_n$：$n$ 個のビンを持つ確率単体、すなわち $\RR^n_+$ における確率ベクトルの集合。
	\item $(\a,\b)$：単体 $\simplex_n \times \simplex_m$ におけるヒストグラム。
	\item $(\al,\be)$：空間 $(\X,\Y)$ 上で定義された測度。
	\item $\frac{\d\al}{\d\be}$：測度 $\al$ の $\be$ に関する相対密度。
	\item $\density{\al} = \frac{\d\al}{\d x}$：測度 $\al$ の Lebesgue 測度に関する密度。
	\item $(\al=\sum_i \a_i \delta_{x_i},\be=\sum_j \b_j \delta_{y_j})$：$x_1,\dots,x_n\in\X$ および $y_1,\dots,y_m\in\Y$ 上に台を持つ離散測度。
	\item $\c(x,y)$：ground cost。対応する対ごとのコスト行列は $\C_{i,j}=(\c(x_i,y_j))_{i,j}$ であり、$\al,\be$ の台上で評価される。
	\item $\pi$：$\al$ と $\be$ の間のカップリング測度。すなわち、任意の $A\subset\X$ に対し $\pi(A\times \Y)= \al(A)$、任意の部分集合 $B\subset \Y$ に対し $\pi(\X\times B)= \be(B)$。離散測度の場合 $\pi = \sum_{i,j} \P_{i,j}\delta_{(x_i,y_j)}$。
	\item $\Couplings(\al,\be)$：カップリング測度の集合。離散測度の場合は $\CouplingsD(\a,\b)$。
	\item $\Potentials(\c)$：許容される双対ポテンシャルの集合。離散測度の場合は $\PotentialsD(\C)$。
	\item $\T : \X \rightarrow \Y$：Monge 写像。通常 $\T_\sharp \al = \be$ を満たす。
	\item $(\al_t)_{t=0}^1$：動的測度。$\al_{t=0}=\al_0$ かつ $\al_{t=1}=\al_1$。
	\item $\speed$：Benamou--Brenier 定式化における速度。$\Moment=\al \speed$：運動量。
	\item $(\f,\g)$：双対ポテンシャル。離散測度の場合 $(\fD,\gD)$ は双対変数。
	\item $(\uD,\vD) \eqdef (e^{\fD/\varepsilon},e^{\gD/\varepsilon})$：Sinkhorn スケーリング。
	\item $\K \eqdef e^{-\C/\varepsilon}$：Sinkhorn のための Gibbs カーネル。
	\item $\flow$：$\Wass_1$ 型問題のためのフロー（発散制約下の最適化）。
	\item $\MKD_\C(\a,\b)$ および $\MK_\c(\al,\be)$：コスト $\C$（ヒストグラム）および $\c$（任意の測度）を持つ OT に関連する最適化問題の値。
	\item $\WassD_p(\a,\b)$ および $\Wass_p(\al,\be)$：ground 距離行列 $\distD$（ヒストグラム）および距離 $\dist$（任意の測度）に関連する $p$-Wasserstein 距離。
	\item $\lambda\in\simplex_S$：$S$ 個の測度の重心を計算するために用いる重みベクトル。
	\item $\dotp{\cdot}{\cdot}$：ベクトル間の通常のユークリッド内積。同サイズの2つの行列 $A$ と $B$ に対し、$\dotp{A}{B} \eqdef \trace(A^\top B)$ は Frobenius 内積。
	\item $\f\oplus \g(x,y) \eqdef f(x)+g(y)$：2つの関数 $\f : \X \rightarrow \RR, \g : \Y \rightarrow \RR$ に対し、$\f \oplus \g : \X\times \Y \rightarrow \RR$ を定義する。
	\item $\fD \oplus \gD \eqdef \fD \ones_m^\top + \ones_n \gD^\top \in \RR^{n\times m}$：2つのベクトル $\fD\in\RR^n$, $\gD\in\RR^m$ に対する定義。
	\item $\al \otimes \be$ は $\X \times \Y$ 上の積測度。\ie
		$\int_{\X \times \Y} g(x,y) \d(\al\otimes\be)(x,y) \eqdef \int_{\X \times \Y} g(x,y) \d\al(x) \d\be(y)$。
	\item $\a \otimes \b \eqdef \a \b^\top \in \RR^{n \times m}$。
	\item $\uD \odot \vD = (\uD_i \vD_i) \in \RR^n$：$(\uD, \vD) \in (\RR^n)^2$ に対する要素ごとの積。
\end{itemize}
